\textbf{1. Rotating Cylinder}

(a) [1] Alice shows Bob her new cylinder collection. One of them is a hollow
cylinder. The cylinder has length $L$, radius $r$ and wall thickness
$\delta \ll r$ for both the bases and the curved surface. The mass density of
the cylinder is $\rho$. Find the mass $m$ of the cylinder and the moment of
inertia $I$ of the cylinder about its axis.

(b) [2] The cylinder is made of carbon nanotubes and has high tensile strength,
$\sigma = \frac{F_{max}}{A}$, where $F_{max}$ is the maximum tension force the
cylinder can sustain without breaking and $A$ is the cross-sectional area
perpendicular to the direction the force is applied (see Fig.\ 1-1). By
considering an element of the cylinder when spinning (refer to the top view of
the spinning cylinder in Fig 1-2), estimate the maximum angular velocity
$\omega_{max}$ at which the cylinder can spin before breaking. Express in terms
of $\sigma, \rho$ and $r$. Give a numerical value of $\omega_{max}$ for
$L = 1m$, $r = 10cm$, $\delta = 50\mu m$, $\sigma = 20GPa$,
$\rho = 1.3g/cm^3$. Ignore deformation. You may assume the sides break before
the bases do.

\begin{center}
\begin{tabular}{ccc}
\includegraphics[width=0.10\linewidth]{figures/PanPhO_2026_Q1_fig1.png} &
\includegraphics[width=0.18\linewidth]{figures/PanPhO_2026_Q1_fig2.png} &
\includegraphics[width=0.30\linewidth]{figures/PanPhO_2026_Q1_fig3.png} \\
 & Top view & \\
Fig 1-1 & Fig 1-2 & Fig 1-3
\end{tabular}
\end{center}

In what follows, we shall assume $\omega \ll \omega_{max}$.

(c) [2] In Fig 1-3, Alice uses a motor positioned at a height $h$ above the
ground to set the cylinder spinning with an initial angular velocity
$\omega_0$. In a prank, Bob uses his \textquotedbl psychic powers\textquotedbl{} to release the cylinder
from the motor without altering its angular velocity. The cylinder, starting
with a negligible initial translational velocity, falls toward the ground,
which has a coefficient of friction $\mu$. Ignoring air resistance, assume that
\underline{collisions} \underline{with} \underline{the} \underline{ground}
\underline{are} \underline{perfectly} \underline{elastic} \underline{in}
\underline{the} \underline{$y$-direction} and that the collision duration
$t_c$ is negligible ($gt_c^2 \ll h$), where $g$ represents the acceleration
due to gravity. Sketch a qualitative graph representing the trajectory of the
cylinder's center of mass.

(d) The cylinder hits the ground at $O$ at $t = 0$.

(di) [2] Find the velocity $(v_x, v_y)$ and the angular velocity $\omega$ of
the cylinder right after the 1\textsuperscript{st} collision.

(dii) [3] Find the position of the ball $x(t)$ and $y(t)$, as well as the
angular velocity $\omega(t)$. Express your answers in terms of the number of
collisions with the ground $N$ (We take $N = 1$ at $t = 0$), the time elapsed
since the last collision $\Delta t, h, g, r,$ and the ratio
$K = \frac{mr^2}{I}$. You can use the floor function $\lfloor x \rfloor$ which
denotes the integer part of $x$ (e.g.\ $\lfloor \pi \rfloor = 3, \lfloor e
\rfloor = 2$).
